% 本文件可直接置入第三章“仿真与结果分析”节。主文件加载 graphicx、booktabs。
% 从模板目录编译时，将 \ChThreeFigurePath 定义为相对主文件的 figures 路径。
\providecommand{\ChThreeFigurePath}{figures}
\providecommand{\ChThreeFigureWidth}{\linewidth}
\providecommand{\ChThreeGroupBreak}{}
\providecommand{\ChThreeMainBreak}{}
\providecommand{\bicaption}[2]{\caption{#1}}
\begingroup
\newcommand{\ChThreeFigure}[4]{%
  \begin{figure}
  \centering
  \includegraphics[width=#4\ChThreeFigureWidth]{\ChThreeFigurePath/fig#1}
  \bicaption{#2}{#3}\label{fig:ch3formal-#1}
  \end{figure}}
\subsection{实验设置与评价指标}
在冻结的合成异构平台上，四个场景共处理500个请求，每个请求的相对期限为40\,ms。CPU、GPU和共享DLA频率各取三个档位，各场景有27种频率组合。CPU、GPU、DLA0和DLA1的初温均为30\,°C。全窗系统能耗计入10.00\,s服务期及末尾40\,ms；所有请求按期完成且四节点连续峰温上界均不超过80\,°C时，方案记为安全（SAFE）。

常温主实验在25\,°C下交叉组合场景正序、逆序和两种驻留时长，每个工况开放全部81个代表部署。正序原驻留的四段帧数依次为62、63、62、63，交替驻留为75、50、75、50，逆序相应反转。Kim-Adapt及其安全增强版均采用24个体、64代，增强版按同一候选排序回退至首个安全方案。热感知遗传算法（Thermal-GA）采用32个体、最多96代。每个随机方法每工况运行5个种子，单次墙钟上限120\,s，两版Kim使用配对种子。Kim在95\,s时截止生成，连续认证和GA搜索在116\,s时截止接纳结果，达到代数上限可提前结束。

记同一实例、同一候选域的最低安全能耗为$E^\star$。本文连续热联合优化方法（CT-Joint）在81部署域的参照值取自M2历史精确证书，四部署参照取相应域的归档证书。本章比较安全性和方案质量，M2的求解加速效果在第四章分析。仅对安全解计算最优差距$100(E-E^\star)/E^\star$；比较相对某基线的节能率时，以该次基线的安全能耗$E$为分母，逐种子计算后取均值。

\subsection{常温下的算法综合效果}
四个工况中，三个随机方法各有20/20次返回安全方案。表~\ref{tab:ch3formal-main}给出安全能耗均值，图~\ref{fig:ch3formal-A}保留逐种子分布。Kim两版在全部20对运行中均由首选通过连续认证，因此两版能耗相同。两个正序工况下，Kim的平均能耗仅比参照高0.000194和0.000182\,J；两个逆序工况下，五个种子均达到参照值。

\begin{table}\centering
\bicaption{25\,°C下四个工况的安全系统能耗}{Safe system energy in four scenarios at 25\,°C}
\label{tab:ch3formal-main}
\small\setlength{\tabcolsep}{4pt}
\begin{tabular}{lrrrr}\toprule
工况 & 本文方法（J） & Kim两版各自（J） & Thermal-GA（J） & GA最优差距（\%）\\\midrule
正序原驻留   & 43.110688 & 43.110883 & 43.241507 & 0.303\\
正序交替驻留 & 43.207258 & 43.207440 & 43.429815 & 0.515\\
逆序原驻留   & 43.250300 & 43.250300 & 44.131804 & 2.038\\
逆序交替驻留 & 43.351353 & 43.351353 & 44.739794 & 3.203\\\bottomrule
\end{tabular}
\end{table}
\ChThreeMainBreak
\ChThreeFigure{A}{25\,°C、81部署下各方法的安全能耗与同域最优差距}{Safe energy and optimality gaps over 81 deployments at 25\,°C}{1}

图~\ref{fig:ch3formal-B}以未舍入能耗统计命中次数，判据为SAFE且$|E-E^\star|\leq 10^{-9}$\,J。每个随机方法的分母均为五次运行，确定性证书只提供$E^\star$。Kim两版各命中10/20次，Thermal-GA命中3/20次。尽管GA在两个正序工况分别命中2次和1次，其余种子的能耗抬高了均值。相对GA，本文方法的平均节能率依次为0.301\%、0.508\%、1.973\%和3.088\%。GA同时搜索部署与四段频率，逆序交替驻留下的单次最优差距最高达4.812\%；这五个种子说明了给定预算下的波动范围。
\ChThreeFigure{B}{25\,°C下随机方法达到81部署域最低安全能耗的次数（命中数/5）}{Runs attaining the 81-deployment minimum safe energy at 25\,°C (hits per five seeds)}{1}

\ChThreeGroupBreak
\subsection{跨场景连续热状态对安全判定的影响}
在25\,°C正序原驻留工况，固定部署C并在四个场景均使用最高频率组合HHH。场景独立评价将各段入口温度重置为30\,°C，连续评价则继承上一段末温。相同方案在两种评价下的峰温上界分别为77.553\,°C和80.773\,°C，前者判安全、后者认证超温。图~\ref{fig:ch3formal-C}的GPU轨迹显示，差异从第一次场景切换后出现。保持部署和负载，改用LLL/LLL/LHL/LHL后，连续峰温上界降至76.506\,°C，系统能耗为43.110758\,J，最小期限裕量为2.721\,ms。该方案的四节点记录含3\,504行事件采样，覆盖末尾计能窗。
\ChThreeFigure{C}{部署C的GPU温度轨迹（图例数值为认证峰温上界）}{GPU temperature under deployment C with certified peak upper bounds in the legends}{1}

误收率以重置评价判安全的频率序列数为分母。表~\ref{tab:ch3formal-reset}和图~\ref{fig:ch3formal-D}(a)表明，部署A和C合计有3.99\%的重置接纳方案在连续运行时超温。图~\ref{fig:ch3formal-D}(b)统计既有27.5\,°C工况中满足期限的序列有多少通过连续热认证。在三个已有分类的场景配置中，A和C各自的59\,049条按期序列均超温，D001和D015的按期序列则全部安全。这些数值描述部署的候选域；算法高温返回率取自后文逆序原驻留的五种子实验。
\begin{table}\centering
\bicaption{25\,°C正序原驻留下逐场景重置的误收情况}{False acceptance by scene-wise reset at 25\,°C in forward order with the original dwell times}
\label{tab:ch3formal-reset}
\small
\begin{tabular}{lrrrr}\toprule
部署 & 重置判安全数 & 其中连续安全数 & 误收数 & 误收率（\%）\\\midrule
A & 59\,049 & 56\,307 & 2\,742 & 4.64\\
C & 59\,049 & 57\,084 & 1\,965 & 3.33\\
合计 & 118\,098 & 113\,391 & 4\,707 & 3.99\\\bottomrule
\end{tabular}
\end{table}
\ChThreeFigure{D}{25\,°C下的重置误收率与27.5\,°C四部署域的连续热分类}{False acceptance under reset evaluation at 25\,°C and continuous thermal classification at 27.5\,°C}{1}

\ChThreeGroupBreak
\subsection{部署选择与逐场景调频的消融分析}
在25\,°C逆序原驻留下，分别改变可选部署和可选频率策略。固定部署使用历史对照A，整窗同频要求四个场景使用同一组CPU、GPU和共享DLA频率。图~\ref{fig:ch3formal-E}和表~\ref{tab:ch3formal-ablation}均用各受限域的最低安全能耗，节能率以固定A、整窗同频的49.029417\,J为基准。
\begin{table}\centering
\bicaption{部署选择与场景调频的四种组合}{Ablation of deployment selection and per-scene frequency control}
\label{tab:ch3formal-ablation}
\small
\begin{tabular}{llrrr}\toprule
部署范围 & 频率策略 & 最优部署 & 能耗（J） & 节能率（\%）\\\midrule
固定A & 整窗同频 & A & 49.029417 & 0.000\\
固定A & 逐场景调频 & A & 43.303623 & 11.678\\
81个可选 & 整窗同频 & D000 & 48.279859 & 1.529\\
81个可选 & 逐场景调频 & D000 & 43.250300 & 11.787\\\bottomrule
\end{tabular}
\end{table}
\ChThreeFigure{E}{四种部署与调频策略的最低安全能耗及相对节能率}{Minimum safe energy and relative savings under four deployment and frequency policies}{1}

固定A时，逐场景调频节能11.678\%；保持整窗同频、只开放部署选择时，节能1.529\%。两项决策均开放时节能11.787\%。从固定A的逐场景方案再开放部署选择，能耗由43.303623降至43.250300\,J，额外下降0.123\%，该比例以43.303623\,J为分母。在本工况81个代表中，68个在两种频率策略下均有安全方案，逐一固定后的调频节能率为9.84\%--18.62\%，中位数17.22\%；另13个在固定部署下均无安全方案。因此，固定A的四格对照主要体现调频的能耗贡献，跨部署选择还影响可行性。

\ChThreeGroupBreak
\subsection{共同四部署域下的热压力比较}
温度对照固定逆序原驻留、同一500请求流及A、C、D001、D015四个部署，两侧分别使用25\,°C和27.5\,°C。图~\ref{fig:ch3formal-F}给出安全返回数。表~\ref{tab:ch3formal-pressure}对随机方法返回的安全方案取能耗和热裕量的算术均值，图~\ref{fig:ch3formal-G}保留逐次结果，热裕量为80\,°C减去认证峰温上界。25\,°C时，两版Kim在五个种子中均首选部署C并达到域内最低值43.250327\,J。
\ChThreeFigure{F}{两种环境温度下的安全返回数与有限域参照}{Safe returns and the finite-domain reference at two ambient temperatures}{.93}

\begin{table}\centering
\bicaption{逆序原驻留下共同四部署域的压力对照}{Matched thermal-pressure comparison for four deployments in reverse order with the original dwell times}
\label{tab:ch3formal-pressure}
\small\setlength{\tabcolsep}{4pt}
\begin{tabular}{lcrcrr}\toprule
& \multicolumn{2}{c}{25\,°C} & \multicolumn{3}{c}{27.5\,°C}\\
\cmidrule(lr){2-3}\cmidrule(lr){4-6}
方法 & 安全 & 能耗（J） & 安全 & 能耗（J） & 热裕量（°C）\\\midrule
本文方法 & 已证 & 43.250327 & 已证 & 45.100779 & 3.680\\
Kim-Adapt & 5/5 & 43.250327 & 0/5 & --- & ---\\
Kim-Adapt+SAFE & 5/5 & 43.250327 & 5/5 & 45.276228 & 9.355\\
Thermal-GA & 5/5 & 43.344397 & 5/5 & 45.206048 & 7.085\\\bottomrule
\end{tabular}
\end{table}

升至27.5\,°C后，Kim首选C仍完成全部请求，实耗43.547995\,J，峰温达81.922\,°C。五次均因超温而排除出安全能耗统计。增强版五次均依次排除C、A，经两次回退选中D015。该方案的热裕量为9.355\,°C，能耗为45.276228\,J。本文方法选中D001，以3.680\,°C热裕量和0.309\,ms最小期限裕量满足约束，比增强版节省0.175448\,J（0.388\%）。D015则多出5.675\,°C热裕量，反映首个安全候选与最低能耗候选之间的取舍。
\ChThreeFigure{G}{共同四部署域中安全方案的系统能耗与热裕量}{System energy and thermal headroom of safe solutions in the matched four-deployment domain}{1}

Thermal-GA在高温下五次均安全，其中两次找到D001的最低能耗方案，三次得到与增强版相同的D015方案。升温使四部署域的最低安全能耗从43.250327升至45.100779\,J，最优部署也由C变为D001。该温度对照仅覆盖冻结合成平台的一种场景次序和驻留配置，每组五个随机种子。
\endgroup
